% !TeX root = BTechProjectTemplate.tex
\chapter*{ABSTRACT}
\addcontentsline{toc}{chapter}{ABSTRACT}
High-resolution Magnetic Resonance Imaging (MRI) assumes an important place in medical neuro-diagnosis as it offers a non-invasive method of observing the complex anatomy of the brain in high resolution. However, there is a catch in getting the latter. It involves some physical laws that create inevitable trade-offs in the resolution, SNR, and acquisition time of the imaging procedure that can only be achieved with the help of the corresponding equipment and not just hardware alone. Multi-modal MRIs such as the T1, T2, and PD MRIs offer rich complimentary data on the anatomy of the brain, but integrating them into the process of image enhancement proves difficult due to their non-linearity. Furthermore, the local receptive field of CNNs is ill-equipped for analyzing global brain anatomies.

In order to handle these issues, we introduce \textbf{Swin-MMHCA}, an end-to-end multiple task framework where a hierarchical Swin-Transformer-V2 encoder is complemented with the Multi-Modal Hybrid Cross Attention (MMHCA) mechanism. The former exploits multi-scale information in a way that still considers long-range dependencies. Since these characteristics are mandatory to ensure anatomical fidelity, an additional structural edge guidance branch is considered to explicitly incorporate high frequency information from the edges, and this is fused with the extracted hierarchical features thanks to the attention mechanism. As a result, sharp anatomical details contained in the target modality (T1, PD) are located and transferred into T2, without being blended via spatial invariant operations.

This constraint is even more emphasized via the segmentation and lesion detection auxiliary tasks that regularize the network, thus forcing the representation learning process not to generate metrically accurate yet anatomically inconsistent structures near tissue edges.
With respect to the IXI data set, our approach obtains a PSNR of 40.62~dB and SSIM of 0.9927 at the $2\times$ upsampling scale while still maintaining a PSNR of 33.32~dB even under the harder $4\times$ magnification task – surpassing all baseline CNN approaches and unimodal models. This study not only demonstrates state-of-the-art performance on commonly used evaluation metrics but also highlights the potential for application in practice, where our approach can be used to enhance clinical acquisition sequences of relatively lower resolution to satisfy an otherwise higher image acquisition time requirement.